\section{Component holonomy and neck degeneration}\label{sec:component-holonomy}

We next allow the component connections to vary.
The compact family in this section has one base
parameter and four necks of length $L+O(1)$.
The bundle parameters are independent of $L$.
The conclusions therefore concern arbitrary
component-to-neck scale ratios; they do not
combine varying component holonomy with the
six independent node parameters of
Section~\ref{sec:geometry}.

\subsection{A compact family and its flat bundles}
Let $E$ be the projective elliptic curve
$y^2=(x-1)(x-2)(x-3)$, with origin $o$ at infinity.
The sections $s_0=x$ and $s_1=1$ of
$\mathcal O_E(2o)$ have disjoint zero divisors
$p+q$ and $2o$. In $E\times\PP^1\times\Delta$ set
\begin{equation}\label{eq:elliptic-family}
 \cC=\{s_0uv+t s_1(u^2+v^2)=0\},
 \qquad \pi:\cC\longrightarrow\Delta .
\end{equation}
Give $E$ a fixed flat Kähler metric and $\PP^1$
a fixed smooth Kähler metric. Restrict their
product, with a positive disc term, to $\cC$.

\begin{proposition}\label{prop:elliptic-family}
After shrinking $\Delta$, \eqref{eq:elliptic-family}
is a proper projective semistable family with
smooth total surface and connected genus-three
smooth fibres. Its central components are two
copies $E_+,E_-$ of $E$ and two rational curves
$R_p,R_q$, with dual graph $K_{2,2}$.

Fix trace-free Hermitian matrices $T_i$ and
nonzero real harmonic one-forms $\omega_i$ on
$E_i$, for $i=+,-$. The component connections
\[
 D_i=d-i\alpha_iT_i\omega_i,\qquad \alpha_i\in\RR,
\]
the trivial connections on $R_p,R_q$, and arbitrary
node identifications $U_{ij}\in\SU(r)$ define
flat determinant-one bundles on a neighborhood
of the whole central fibre. For every fixed
$\alpha$ they are holomorphic bundles over the
restricted family, with parallel metric $h_0$.
Dependence on $\alpha$ is smooth and real.
\end{proposition}
\begin{proof}
At a node the equation is $xw+t(1+w^2)=0$.
Its derivative in $t$ is nonzero there.
Elsewhere on the central fibre a spatial
derivative is nonzero, including above $o$:
divide by $s_0$ and use that $s_1/s_0$ has a
double zero. Compactness gives a smooth total
surface after restriction. The projection is
proper and projective. Its fibres are effective
Cartier divisors, and the local map to the disc
is flat. The holomorphic Morse normal form
gives $ab=t$ at each node.

For $t\ne0$, projection to $E$ has degree two.
The discriminant in the affine chart is
$x^2-4t^2$, with four distinct simple zeros
for small $t$. At $o$ the equation
$w+t(s_1/s_0)(1+w^2)=0$ and its reciprocal
chart give two unramified points.
Thus $C_t$ is smooth. A disconnected
double cover would be unramified, so it is
connected. Riemann--Hurwitz gives genus three.

Use disjoint node bidiscs and component
exteriors as in Proposition~\ref{prop:bundle}.
Trivialize each exterior smoothly over the
disc, with the lifts agreeing with the branch
coordinate on its collars. Pull back
$\omega_i$ from the fixed punctured component.
It is closed and has zero puncture periods.
On an incidence collar write $\omega_i=dF_{ie}$,
with the central primitive normalized to
vanish at the node. The frame change
$G_{ie}=\exp(i\alpha_iT_iF_{ie})$ gives a
parallel unitary frame. Give each node bidisc
the trivial flat bundle and identify these
parallel frames by the prescribed constant
unitary maps. Only incidence overlaps occur;
there are no triple overlaps.

These are flat isomorphisms, so the connections
glue to a unitary flat connection with parallel
determinant. Its $(0,1)$ part is integrable.
The exterior smooth frames need not be
holomorphic; the local parallel frames are.
Properness permits a smaller disc on which
this neighborhood contains every whole fibre.
This proves the claimed total-family construction.
\end{proof}

Let $A_i$ be the four limiting areas and put
$\kappa_i=\int_{E_i}|\omega_i|^2$.
The elliptic limiting metrics are flat, so
$|\omega_i|^2=\kappa_i/A_i$ is constant.
All necks have leading capacity $2\pi/L$, where
$L=\log(1/|t|)$.

The distinction between node identifications
and based holonomy matters here.
Killing the annular core loops gives a
surjection
\[
 \pi_1(C_t)\longrightarrow
 \pi_1(E_+)*\pi_1(E_-)*\mathbb Z .
\]
Fix paths $c_{ie}$ from component basepoints
to the nodes and put
\[
 \mathcal U_{ie}=\exp\left(i\alpha_iT_i
                       \int_{c_{ie}}\omega_i\right),
 \qquad K_{ij}=\mathcal U_{ie}^{-1}U_{ij}\mathcal U_{je}.
\]
On rational components $\mathcal U_{ie}=I$.
Spanning-tree products of the $K_{ij}$ identify
the component basepoint fibres, and their cycle
product represents the chosen based graph loop.
A bare product of the $U_{ij}$ need not do so.
Changing paths changes the presentation of
the same local system.

For all sufficiently small $\alpha_i$, the
parallel trace-free Hermitian endomorphisms
are exactly the tuples satisfying
\begin{equation}\label{eq:coupled-kernel}
 X_i=\Ad(U_{ij})X_j,\qquad
 \alpha_i[T_i,X_i]=0\quad(i=+,-).
\end{equation}
Indeed, choose a nonzero period of $\omega_i$.
For small nonzero $\alpha_i$, exponentiation
does not identify distinct eigenvalues of
$T_i$ at this period. Its holonomy commutant
is therefore the commutant of $T_i$.
Such an endomorphism is constant in the
harmonic component gauge and unchanged by
each $G_{ie}$; node compatibility is precisely
\eqref{eq:coupled-kernel}. At $\alpha_i=0$
the component condition is vacuous.
The unitary correspondence \cite{narasimhan-seshadri} gives
polystability, and stability exactly when this
trace-free commutant is zero.

For example, take $r=2$, $T_+=\sigma_z$,
$\alpha_-=0$, and a harmonic form on $E_+$
with periods $(1,0)$. Set the tree maps to
$I$ and the remaining node map to
$W=(I+i\sigma_x)/\sqrt2$.
Tree compatibility makes all $X_i=X$.
When $\alpha_+\ne0$ is small the component
condition gives $X=c\sigma_z$, and
$W\sigma_zW^*=\sigma_y$ forces $c=0$.
The bundle is stable. At $\alpha_+=0$ the
representation is generated by $W$ and is
reducible. This is actual stable-to-polystable
variation, independent of $t$. It does not
produce an unstable locus or a topologically
nontrivial smooth $\SU(r)$ bundle on a curve.

\subsection{The coupled relative spectrum}
On $\cV=\bigoplus_{i=1}^4\Herm_0(r)$ define
\begin{equation}\label{eq:coupled-form}
 \begin{split}
 g_\alpha(X)&=\sum_{i=+,-}\alpha_i^2\kappa_i
                                  \|[T_i,X_i]\|_{\HS}^2,\\
 q_U(X)&=\sum_{e=ij}\|X_i-\Ad(U_{ij})X_j\|_{\HS}^2,\\
 Q_{t,\alpha,U}(X)&=g_\alpha(X)+\frac{2\pi}{L}q_U(X),
 \qquad M(X)=\sum_i A_i\tr(X_i^2).
 \end{split}
\end{equation}
Write $\alpha_*=\max_i|\alpha_i|$,
$m=4(r^2-1)$, and let $\nu_0,\ldots,\nu_{m-1}$
be the generalized eigenvalues of $(Q,M)$,
including zeros.

\begin{theorem}\label{thm:coupled-spectrum}
Let $\lambda_k$ be the eigenvalues of $D^*D$
on the real trace-free Hermitian endomorphisms
of the bundle in Proposition~\ref{prop:elliptic-family}.
As $(|t|,\alpha_*)\to(0,0)$,
\[
 (1-\epsilon)\nu_k\le\lambda_k\le(1+\epsilon)\nu_k
 \quad(0\le k<m),\qquad \epsilon\longrightarrow0,
\]
uniformly in all $U$. There is a uniform
complementary gap $\lambda_m\ge\gamma>0$
in a sufficiently small product neighborhood.
The exact kernel is \eqref{eq:coupled-kernel}.
No restriction is imposed on $\alpha_i^2L$
or on the distance to a reducible representation.
\end{theorem}
\begin{proof}
We give the component and matching estimates
that replace those of Section~\ref{sec:spectrum}.
On a closed elliptic component a constant
matrix has energy
$\alpha_i^2\kappa_i\|[T_i,X_i]\|^2$.
In flat coordinates the constant subspace is
invariant under $D_i^*D_i$. After diagonalizing
$T_i$, entry $ab$ on a dual-lattice Fourier
mode $k$ has eigenvalue
$|2\pi k-\alpha_i(\lambda_a(T_i)-\lambda_b(T_i))
\omega_i|^2$. The nonzero Fourier modes have
a uniform gap for small $\alpha_i$.

Fix a small radius $a>0$. A local primitive
$F$ satisfies $F(0)=0$, $|F|\le Ca$ and
$|dF|\le C$. On $a<|z|<2a$ join $X_i$ to
the parallel section $\Ad(e^{i\alpha_iT_iF})X_i$
with a radial cutoff. Since
\[
 \|\Ad(e^{i\alpha_iT_iF})X_i-X_i\|
 \le|\alpha_iF|\,\|[T_i,X_i]\|,
\]
this modification changes the component
energy by at most $Ca^2\alpha_i^2
\|[T_i,X_i]\|^2$. In the flat node frames
interpolate affinely in log radius over
$\ell_a=L+2\log a$.
For the resulting trial section $JX$,
\[
 \begin{split}
 E(JX)&\le(1+Ca^2+o_a(t))g_\alpha(X)
                         +\frac{2\pi}{\ell_a}q_U(X),\\
 |\|JX\|^2-M(X)|&\le(Ca^2+o_a(t))M(X).
 \end{split}
\]
Metric transfer errors are multiplicative
because the pulled-back connection is fixed.
These bounds give the relative upper
estimate and $\lambda_{m-1}\le
C(\alpha_*^2+L^{-1})$.

For the lower estimate, let $\Omega_a$ be a
punctured elliptic component and let $\omega_a$
be the absolute harmonic representative of
the class of $\omega|_{\Omega_a}$.
Thus $\omega_a$ is closed, coclosed, has zero
normal component, and
$\omega=\omega_a+df_a$. Write
$\kappa_a=\|\omega_a\|^2$. Then
$\kappa_a\to\kappa$ as $a\to0$.
Indeed $\omega$ is a competitor, so
$\kappa_a\le\kappa$. A bounded-energy weak
limit retains its periods and is harmonic
away from the punctures. An isolated
singularity of an $L^2$ harmonic one-form
on a surface is removable: even a simple
pole in its holomorphic Laurent expansion
has infinite $L^2$ norm. The limit is
therefore $\omega$, and lower
semicontinuity proves the assertion.

Gauge by $e^{i\alpha T f_a}$. For an entry
with eigenvalue difference $\delta$ of $T$,
the connection becomes $d-i\beta\omega_a$,
where $\beta=\alpha\delta$.
Write $v=z+w$, with $\int_{\Omega_a}w=0$.
The absolute boundary condition gives
$\langle dw,\omega_a z\rangle=0$.
For small $|\beta|$, Poincaré implies
$\|w\|\le C_a\|(d-i\beta\omega_a)w\|$.
The remaining cross term is bounded by
$C_a\beta^2|z|\|w\|$. Completing the square
therefore gives
\[
 \begin{split}
 \|(d-i\beta\omega_a)v\|^2
 &\ge\tfrac12\|(d-i\beta\omega_a)w\|^2
                +(\beta^2\kappa_a-C_a\beta^4)|z|^2,\\
 \|w\|_{H^1}^2+\beta^2|z|^2
 &\le C_a\|(d-i\beta\omega_a)v\|^2 .
 \end{split}
\]
These inequalities remain valid at $\beta=0$.
Summing entries defines component means $X_i$
and gives, with multiplicative $1+o_a(t)$
changes on the nearby fibre,
\[
 \begin{split}
 E_i&\ge(1-C_a\alpha_*^2)
                  \alpha_i^2\kappa_{i,a}\|[T_i,X_i]\|^2,\\
 \|w_i\|_{H^1}^2+\alpha_i^2\|[T_i,X_i]\|^2
 &\le C_a E_i.
 \end{split}
\]
On rational components take the ordinary mean.

The mean gauge and the parallel node frame
differ by $e^{i\alpha_iT_i(f_{i,a}-F_{ie})}$.
On a fixed collar its action on $X_i$ differs
from $X_i$ by at most
$C_a|\alpha_i|\|[T_i,X_i]\|$.
The trace bound on $w_i$ thus gives
$|c_{ie}-X_i|\le C_a E_i^{1/2}$ for the
parallel-frame boundary averages.
In particular there is no additive error
$O(\alpha_i)\|X_i\|$.
Separate the disjoint component and neck
energies as $E_{\mathrm c}$ and $E_{\mathrm n}$.
Angular averaging on the annuli yields
\[
 q_U(X)^{1/2}
 \le\sqrt{\ell_a/(2\pi)}\,E_{\mathrm n}^{1/2}
                                +C_aE_{\mathrm c}^{1/2}.
 \]
After multiplication by $\sqrt{2\pi/L}$ and
Young's inequality, this gives
\[
 \frac{2\pi}{L}q_U(X)
 \le(1+O_a(L^{-1/2}))E_{\mathrm n}
                              +O_a(L^{-1/2})E_{\mathrm c}.
\]
Combining it with the component estimate,
without counting either energy twice, gives
\begin{equation}\label{eq:coupled-lower}
 Q(X(s))\le
 \bigl(1+\theta(a)+C_a\alpha_*^2
              +O_a(L^{-1/2})+o_a(t)\bigr)E(s),
 \qquad \theta(a)\to0 .
\end{equation}

We also need the mass comparison. For a
normalized eigensection of eigenvalue
$\lambda\le1$, the preceding bounds give
$\|w_i\|_{H^1}\le C_a\sqrt\lambda$ and bounded
means. In the mean gauge the central operator
applied to $X_i$ is
$\alpha_i^2|\omega_{i,a}|^2\operatorname{ad}(T_i)^2X_i$,
bounded by $C_a|\alpha_i|\sqrt{E_i}$.
The metric perturbation adds at most
$o_a(t)|\alpha_i|\|[T_i,X_i]\|$.
Interior $W^{2,2}$ estimates on fixed enlarged
collars therefore control the boundary
values in $C^0$. The enlargement passes
slightly into the flat node chart and
introduces no shrinking-domain estimate.
Harmonic extension in that chart, followed
by Lemma~\ref{lem:zero-boundary} for its
zero-boundary remainder, rules out annular
mass loss. For a combination of at most
$m+1$ eigenvectors of eigenvalues at most
$\Lambda\le1$, the result is
\[
 |M(X(s))-\|s\|^2|
 \le(Ca^2+C_a\sqrt\Lambda+o_a(t))\|s\|^2 .
\]
This is the mass argument of
Lemma~\ref{lem:low-mass} with the
energy-controlled component gauges above.

If $\lambda_m$ tended to zero, these means
would inject an $(m+1)$-dimensional space
into $\cV$, a contradiction. The same mass
estimate and \eqref{eq:coupled-lower},
applied to the first $k+1$ eigensections,
give the relative lower estimate by min--max.
Choose $a$ first and then choose $|t|$ and
$\alpha_*$ independently small. This proves
uniformity for every ratio $\alpha_i^2L$.
Finally $Q=0$ is exactly
\eqref{eq:coupled-kernel}, already identified
with the actual parallel sections.
\end{proof}

\subsection{The exact component field}
For $H\in\cP_r$ set
$f_T(H)=\tr(TH^{-1}TH)-\tr(T^2)$.
In the smooth harmonic component frame a
spatially constant metric has
\[
 H^{-1}\partial_{D_i}H
 =-i\alpha_i(H^{-1}T_iH-T_i)\omega_i^{1,0},
 \qquad
 F_H=\alpha_i^2[T_i,H^{-1}T_iH]\,
                   \omega_i^{1,0}\wedge\omega_i^{0,1}.
\]
Since $2i\Lambda(\omega_i^{1,0}\wedge
\omega_i^{0,1})=|\omega_i|^2$, its exact
component HYM equation is
\[
 H^{-1}\dot H=-\alpha_i^2|\omega_i|^2[T_i,H^{-1}T_iH].
\]

Differentiation gives
$df_T(HS)=\tr([T,H^{-1}TH]S)$.
The gradient is $H$-self-adjoint and trace
free. It is geodesically convex: along
$H(u)=H_0^{1/2}e^{uK}H_0^{1/2}$, put
$B=H_0^{-1/2}TH_0^{1/2}$ and diagonalize
$K$ with eigenvalues $k_a$. Then
\[
 \tr(TH(u)^{-1}TH(u))
       =\sum_{a,b}|B_{ba}|^2e^{u(k_a-k_b)}.
\]
For $H=e^S$, the Hermitian symmetry of $T$
also gives
\[
 f_T(e^S)=2\sum_{a<b}|T_{ab}|^2
                      (\cosh(s_a-s_b)-1)
 \ge\tfrac12\|[T,S]\|^2,\qquad
 f_T(e^{uS})\le u^2f_T(e^S)\quad(0\le u\le1).
\]

Equip $\cP_r^4$ with the masses $A_i$ and
define
\[
 \Phi_{t,\alpha,U}(H)=\sum_{i=+,-}\alpha_i^2\kappa_i f_{T_i}(H_i)
        +\frac{\pi}{L}\sum_{e=ij}d(H_i,U_{ij}H_jU_{ij}^*)^2.
\]
Its negative gradient is
\begin{equation}\label{eq:coupled-flow}
 A_iH_i^{-1}\dot H_i
 =-\alpha_i^2\kappa_i[T_i,H_i^{-1}T_iH_i]
          +\frac{2\pi}{L}\sum_{j:ij}B_{ij},
 \qquad B_{ij}=\log(H_i^{-1}U_{ij}H_jU_{ij}^*),
\end{equation}
with the component term omitted on rational
vertices. All derivatives use physical time.

Convexity gives nonexpansion of this flow.
Radial contraction towards the identity
gives $\Phi_{t,\alpha,U}(e^{uS})\le u^2\Phi_{t,\alpha,U}(e^S)$,
and the distance-decreasing logarithm gives
$\Phi_{t,\alpha,U}(e^S)\ge Q(S)/2$.
Consequently
\begin{equation}\label{eq:coupled-decay}
 R(s):=\tfrac12 M(\log H(s))
 \le R(0)e^{-2\nu_0s}.
\end{equation}
Indeed $R'\le-2\Phi_{t,\alpha,U}\le-2\nu_0R$.
This also confines every trajectory to its
initial product ball and proves all-time
existence, including $\nu_0=0$.
On that ball its speed, and the physical
time derivative of any fixed smooth tuple
function, are bounded by $C\zeta$, where
$\zeta=\alpha_*^2+L^{-1}$.
The speed in neck time need not be bounded.

\subsection{Actual coupled tracking}
Put
\[
 e(t,\alpha)=\alpha_*+\frac{(1+\log L)^2}{L}.
\]
The base, component and gluing parameters
remain fixed during each trajectory.

\begin{theorem}\label{thm:coupled-tracking}
For $|t|$ and $\alpha_*$ independently
sufficiently small, and initial tuples in a fixed bounded ball,
the finite flow \eqref{eq:coupled-flow}
defines a smooth determinant-one comparison
metric $\widehat K(s)$ on $C_t$, for all
$s\ge0$, with uniform positive bounds and
\begin{equation}\label{eq:coupled-residual}
 \sup_{s\ge0}\int_{C_t}|\mathcal R(\widehat K(s))|_{\op}
 \le C\left(\frac{\alpha_*}{L}
                       +\frac{(1+\log L)^2}{L^2}\right).
\end{equation}
On each fixed compact component region,
in the harmonic gauge,
$d(\widehat K(s),H_i(s))\le C_K/L$.
For an actual HYM flow with fixed positive
initial bounds, let
$\eta=\int_{C_t}\sigma(H(0),\widehat K(0))$.
Then for every $T\ge0$,
\begin{equation}\label{eq:coupled-tracking}
 \sup_{0\le s\le TL}\int_{C_t}d(H(s),\widehat K(s))^2
                       \le\eta+CTe(t,\alpha).
\end{equation}
All constants are independent of $T$, $U$
and the ratios $\alpha_i^2L$.
If in addition $L\nu_0\ge\nu_*>0$, then
\begin{equation}\label{eq:coupled-alltime}
 \sup_{s\ge0}\int_{C_t}d(H(s),\widehat K(s))^2
 \le\eta+C_{\nu_*}e(t,\alpha)(1+|\log e(t,\alpha)|)
\end{equation}
for small $e(t,\alpha)$.
\end{theorem}
\begin{proof}
Use the component Green functions of
Section~\ref{sec:actual-flow} and set
\[
 \xi_i=-\sum_{e=ij}B_{ij}g_{ie},\qquad
 V_i=\frac{2\pi}{A_i}\sum_{j:ij}B_{ij},
 \qquad K_i=H_i e^{\varepsilon\xi_i},
 \quad \varepsilon=L^{-1}.
\]
Off the punctures $\Delta_i\xi_i=-V_i$.
On a component cut at radius $\rho/2$,
the fixed Green estimates give
\[
 \begin{gathered}
 \|\xi_i\|_\infty\le C(1+|\log\rho|),\qquad
 |\nabla^k\xi_i|\le C(1+|u|^{-k})\quad(k=1,2),\\
 \int(|\xi_i|+|\xi_i|^2+|\nabla\xi_i|)\le C,\qquad
 \int|\nabla\xi_i|^2\le C(1+|\log\rho|),\\
 \int|\xi_i|\,|\Delta_i\xi_i|\le C,\qquad
 \int|\dot\xi_i|\le C\zeta .
 \end{gathered}
\]
The derivative bound is written near each
node in its coordinate $u$; away from the
nodes it is uniform. The last estimate uses the physical speed
bound; no bound on $L\dot H_i$ is used.

In the harmonic gauge write $h=\varepsilon\xi_i$
and $C_H=H^{-1}\partial_{D_i}H$.
The exact identity
\[
 K^{-1}\partial_{D_i}K
       =e^{-h}C_He^h+e^{-h}\partial_{D_i}e^h
\]
and the integral formula for the derivative
of the exponential imply, when $|h|\le1$,
\[
 |2i\Lambda F_K-2i\Lambda F_H-\Delta_i h|
 \le C\bigl(
 |\alpha_i||\nabla h|+\alpha_i^2|h|
                +|\nabla h|^2+|h||\Delta_i h|\bigr).
\]
Here $|C_H|=O(|\alpha_i|)$ and its covariant
derivative is $O(\alpha_i^2)$.
The terms linear in $h$ contain either a
connection coefficient and a derivative of
$h$, or two connection coefficients and
$h$. Differentiating the exponential produces
the two quadratic terms. The time identity
$K^{-1}\dot K=e^{-h}(H^{-1}\dot H)e^h+
e^{-h}\partial_s e^h$ has error
$C(|\dot h|+\zeta|h|)$.
Since $|\omega_i|^2=\kappa_i/A_i$,
\eqref{eq:coupled-flow} cancels the leading
terms exactly:
\[
 H_i^{-1}\dot H_i+2i\Lambda_i F_{H_i}
                                +\varepsilon\Delta_i\xi_i=0.
\]
Thus the integrated component residual is
at most
\[
 C\{\varepsilon(\zeta+\alpha_*+\alpha_*^2)
                       +\varepsilon^2(1+|\log\rho|)\}.
\]
There is no uncancelled pure component
error of order $\alpha_*^4$.
Transfer to the smoothing adds
$C|t|\rho^{-N}$ for a fixed $N$.

At incidence $ie$ write
$\xi_i=-B_{ij}\log|u|+c_{ie}+q_{ie}(u)$,
where $q_{ie}(0)=0$, and put
$P_{ie}=H_i e^{\varepsilon c_{ie}}$.
In the node parallel frame use
\[
 N_e=\Geo(P_{ie},U_{ij}P_{je}U_{ij}^*;
                                          -\varepsilon\log|u|).
\]
Geodesic reversal and $uv=t$ give exact
compatibility from both ends.
This node metric has zero spatial Chern
curvature and physical time velocity
$O(\zeta)$, so its middle-node residual
has integral $O(\zeta\rho^2)$.

The component metric in this parallel
frame is $K_i^{\mathrm{par}}=G_{ie}^*K_iG_{ie}$,
where $G_{ie}=e^{i\alpha_iT_iF_{ie}}$ and
$F_{ie}(0)=0$. On the incidence annulus,
in the logarithm $\ell_i(K)=\log(H_i^{-1}K)$,
the endpoint expansion used in
\eqref{eq:mixed-collar} gives
\[
 |D^k(\ell_i(K_i^{\mathrm{par}})-\ell_i(N_e))|
 \le C\{\alpha_*\rho^{1-k}
       +\varepsilon\rho^{1-k}
       +\varepsilon^2(1+|\log\rho|)\rho^{-k}\},
 \quad k=0,1,2.
\]
The new first term comes from $G_{ie}-I$.
Each estimate differentiated in physical
time gains a factor $C\zeta$.
Interpolate the two $H_i$-self-adjoint,
trace-free logarithms and exponentiate.
The curvature-difference estimate of
Section~\ref{sec:actual-flow}, with
$|Dn|\le C\varepsilon/\rho$ and
$|D^2n|\le C\varepsilon/\rho^2$ for
$n=\ell_i(N_e)$, gives the integrated
collar bound
\[
 C\{\alpha_*\rho+\varepsilon\rho
       +\varepsilon^2(1+|\log\rho|)^2
       +\alpha_*^2\rho^2+\zeta\rho^2\}.
\]
Choose $\rho=L^{-1}$.
Then $\varepsilon|\log\rho|\to0$,
the exponentially small transfer error
is absorbed, and all pieces give
\eqref{eq:coupled-residual}.
They remain in a common positive matrix
range for all physical time.
On fixed component regions the Green
correction is $O(L^{-1})$.

Lemma~\ref{lem:trace-comparison} is local
in parallel holomorphic frames and
applies to this flat bundle. Integrating
\eqref{eq:coupled-residual} proves
\eqref{eq:coupled-tracking}, without a
factor exponential in $\alpha_*^2s$.

For the all-time statement,
Theorem~\ref{thm:coupled-spectrum} gives
$\lambda_0\ge\nu_*/(2L)$.
The actual trace coercivity used in
Theorem~\ref{thm:all-time} remains valid
in the local parallel frames:
for $H=e^S$ and $cI\le H\le CI$,
\[
 \frac c2\|S\|_{L^2}^2
 \le F:=\int\tr(H-I)\le\frac C2\|S\|_{L^2}^2,
 \qquad -F'\ge\frac{c^2}{C}\int|DS|^2 .
\]
It implies exponential decay in $s/L$
with rate bounded below by a multiple
of $\nu_*$. Equation~\eqref{eq:coupled-decay}
gives the same type of bound for the
finite flow. All correction coefficients
vanish at the identity tuple, which is
fixed by the component gauge changes.
The reconstructed metric consequently
has an exponential tail as well.
Combine these tails beyond time $TL$
with \eqref{eq:coupled-tracking} before
that time and take $T$ proportional to
$|\log e(t,\alpha)|$. This proves
\eqref{eq:coupled-alltime}.
\end{proof}

The additional rate condition is nonempty
even when $\alpha_*^2L\to\infty$.
For the rank-two example above write
$Q=hq+aP$, where $h=2\pi/L$,
$a=\alpha_+^2\kappa_+$ and
$P(X)=\|[\sigma_z,X_+]\|^2$.
The kernel computation gives
$\ker q\cap\ker P=0$.
If $a/h\ge\eta_0>0$, then
$Q\ge h(q+\eta_0P)$ and $L\nu_0$
has a positive lower bound.
As $a/h\to0$ that bound fails; the
bounded-neck-time theorem still applies.

\subsection{The initial layer with varying component holonomy}
Use the norm \eqref{eq:extension-norm} with
$\rho=L^{-1}$, the four components of
\eqref{eq:elliptic-family}, and the harmonic
component frames. Denote it by
$\mathcal N_{t,\alpha,U}$.
It again defines a fixed-radius nonempty
open class $\mathcal N_{t,\alpha,U}(S)<R_0$
on each smooth fibre, with only a $C^0$
bound imposed in the middle necks.
For admissible component extensions take
the phases $P_i$ of their \emph{trivial}
component flows from
Theorem~\ref{thm:component-phase}.

\begin{theorem}\label{thm:coupled-entry}
In the joint limit $(|t|,\alpha_*)\to(0,0)$,
there is a preparation time
$s_*(t,\alpha;R_0)\to\infty$ such that
$(\alpha_*^2+L^{-1})s_*\to0$ and
\[
 \eta_*=\int_{C_t}\sigma(H(s_*),\widehat K_P(0))
 \longrightarrow0
\]
uniformly over this initial class and
all unitary parameters.
The estimates of
Theorem~\ref{thm:coupled-tracking} apply
after $s_*$ with initial error $\eta_*$.
The finite flow may instead be aligned
with physical time zero at the cost of
a fixed factor in the squared error and
an additional $O((\zeta s_*)^2)$.
\end{theorem}
\begin{proof}
The preparation construction in
Theorem~\ref{thm:entry} uses trivial
component solutions $G_i(s,z)$ with
uniform derivative bounds $C_n$ on
$0\le s\le n$. Such bounds follow
on each fixed component by compactness
of the bounded $C^{6,\beta}$ initial
class in $C^{4,\beta/2}$, continuous
dependence, and finite-time parabolic
regularity. Enlarge $C_n\ge1$ to be
nondecreasing.

In a harmonic frame, replacing ordinary
$\partial$ by $\partial_{D_i}$ introduces
first-order errors $O(C_n|\alpha_i|)$
and zero-order errors $O(C_n\alpha_i^2)$.
In the node frame the additional gauge
factor is $I+O(\alpha_*\rho)$.
Use the spatially constant nodal value
$G_i(s,p_{ie})$ as reference for both
collar logarithms, as in
Theorem~\ref{thm:entry}.
Their derivative difference is bounded
by $C_n(\rho+L^{-1}|\log\rho|)\rho^{-k}$
for $k=0,1,2$.
The middle-neck time error is
$O(C_n\rho^2)$ and transfer costs
$C_n|t|\rho^{-N}$.
Thus the global preparation metric satisfies
\[
 \sup_{0\le s\le n}\int|\mathcal R(K_{\mathrm{pre}})|
 \le C_n\delta,\qquad
 \delta=\alpha_*+\frac{1+\log L}{L}+|t|L^N\to0 .
\]
Its initial discrepancy has integral
$O(\rho^2)$ because the metrics agree
on the exteriors. Initial neck
derivatives do not enter this estimate.

Choose $s_*=\max\{n:nC_n\le\delta^{-1/2}\}$.
Then $s_*\to\infty$,
$s_*C_{s_*}\delta\le\sqrt\delta$ and
$\zeta s_*\le C\sqrt\delta\to0$.
The trace comparison, fixed-component
phase convergence, and the bounded
$L^2$ norms of the Green functions give
\[
 \eta_*\le C_{R_0}
       \{\sqrt\delta+e^{-\gamma_{R_0}(s_*-1)}+L^{-2}\}.
\]
Restart the coupled tracking theorem.

Two admissible extensions differ only
in discs of area $O(\rho^2)$.
The phase contraction estimate gives
an $O(\rho)$ ambiguity in $P$.
Nonexpansion of \eqref{eq:coupled-flow}
preserves this bound for all time.
The reconstruction is uniformly
Lipschitz from the confined tuple ball
to spatial $L^2$: its Green terms have
uniformly bounded $L^2$ coefficients,
and node geodesics and collar
interpolations have bounded endpoint
derivatives. The same observation and
the physical speed bound give
\[
 \int d(H(s_*+v),\widehat K_P(s_*+v))^2
 \le2\int d(H(s_*+v),\widehat K_P(v))^2
                                      +C(\zeta s_*)^2 .
\]
This proves the time alignment and
extension-independence assertions.
\end{proof}

The phase theorem here uses that the
limiting component connections become
trivial as $\alpha\to0$.
It does not select phases for arbitrary
nontrivial limiting local systems.
Uniform all-time comparison retains
$L\nu_0\ge\nu_*>0$; it has not been
deduced at a degenerating smallest rate
from the relative spectrum alone.
